\cxtitle{Strict diagonal dominance does not ensure diminishing Nystr\"om error reductions}

\begin{cxcredits}
\posedby{Mark Fornace \citeyearpar{AmselEtAl2026}}
\foundby{OpenAI Codex}{2026-08}
\formalizedby{OpenAI Codex}
\auditedby{OpenAI Codex}
\contributedby{Suvrit Sra}
\end{cxcredits}

\begin{cxcontext}
For a symmetric positive-definite matrix $L$ and $\gamma>0$, put
$K=(L+\gamma I)^{-1}$.  For a nonempty selected index set $\mathcal I$, the
nuclear Nystr\"om error is
\[
 F(\mathcal I)=\left\|K-K_{:,\mathcal I}
 K_{\mathcal I,\mathcal I}^{-1}K_{\mathcal I,:}\right\|_*.
\]
The residual is positive semidefinite, so its nuclear norm equals its trace.
The intended diminishing-error-reduction property is that, for
$S\subseteq T$ and $i\notin T$,
\[
 F(S)-F(S\cup\{i\})\geq F(T)-F(T\cup\{i\}).
\]
This formulation avoids the convention-dependent choice between the words
``submodular'' and ``supermodular'' for a decreasing error function.
The property is what would license selecting $\mathcal I$ greedily; without it,
landmark sets are chosen by other means, for instance by determinantal sampling
\citep{LiJegelkaSra2016} or through elementary-symmetric-polynomial design
objectives \citep{MarietSra2017}.
Colbrook had already resolved Problem~\ref{ctx:sdd-nystrom-diminishing-returns}
in July 2026, proving the SDDM case
and giving exact SDD obstructions \cite{Colbrook2026}.  The matrix below was
found independently in August 2026; no priority over Colbrook is claimed.
\end{cxcontext}

\begin{cxsource}{sdd-nystrom-diminishing-returns}
Problem~4.6 of Amsel et al. \cite{AmselEtAl2026}, posed there by Mark Fornace.
Colbrook \cite{Colbrook2026} subsequently completed the answer before this
independent witness was found.
\end{cxsource}

\begin{cxstatement}{sdd-nystrom-diminishing-returns}
Prove or disprove the submodularity of the nuclear Nystr\"om error~(14) when
$L$ is assumed, in contrast to the above, to be (a) SDDM and positive-definite
or (b) SDD and positive-definite.
\end{cxstatement}

\begin{cxsummary}{sdd-nystrom-diminishing-returns}
A strictly diagonally dominant positive-definite $4\times4$ signed matrix has
a nonempty-base marginal Nystr\"om error reduction that increases after an
additional index is selected.
\end{cxsummary}

\begin{cxcertificate}{sdd-nystrom-diminishing-returns}
An exact integer $4\times4$ matrix, rational shift, four rational principal
inverse traces, and the positive violation gap $2165/14833896$.
\end{cxcertificate}

\begin{cxrefutation}
Problem~\ref{ctx:sdd-nystrom-diminishing-returns} asks whether the nuclear
Nystr\"om error retains diminishing reductions for positive-definite SDDM or
SDD matrices \cite{AmselEtAl2026}.
Part~(a) is true, while part~(b) is false \cite{Colbrook2026}.  Here is an
independently found exact witness for part~(b) with a nonempty base set.

\begin{theorem}[\statusfalse: SDD Nystr\"om diminishing returns]
\label{thm:sdd-nystrom-diminishing-returns}
The nuclear Nystr\"om error need not have diminishing error reductions when
$L$ is symmetric diagonally dominant and positive definite, even when $L$ is
strictly diagonally dominant and the comparison uses a nonempty base set.
\end{theorem}
\begin{proof}
Let
\[
 M=\begin{pmatrix}
 11&3&-4&3\\
 3&12&4&-1\\
 -4&4&14&-2\\
 3&-1&-2&11
 \end{pmatrix},\qquad
 \gamma=\frac12,\qquad L=M-\frac12I,
\]
and set $K=(L+\gamma I)^{-1}=M^{-1}$.  The diagonal entries of $L$ minus the
sums of the absolute values of the corresponding off-diagonal entries are
\[
 \frac12,\quad\frac72,\quad\frac72,\quad\frac92.
\]
Thus $L$ is symmetric strictly diagonally dominant with positive diagonal;
Gershgorin's theorem gives $L\succ0$.

For a selected set $\mathcal I$ with complement $J$, the block inverse
identity gives
\[
 K_{J,J}-K_{J,\mathcal I}K_{\mathcal I,\mathcal I}^{-1}
 K_{\mathcal I,J}=(M_{J,J})^{-1}.
\]
All other blocks of the Nystr\"om residual vanish.  Consequently
\[
 F(\mathcal I)=\operatorname{Tr}\bigl((M_{J,J})^{-1}\bigr).
\]

Use one-based indices, take the nonempty base $S=\{2\}$, and compare adding
$i=3$ before and after $j=4$.  Direct exact inversion of the four relevant
principal submatrices gives
\[
 F(\{2\})=\frac{100}{349},\qquad
 F(\{2,3\})=\frac{11}{56},\qquad
 F(\{2,4\})=\frac{25}{138},\qquad
 F(\{2,3,4\})=\frac1{11}.
\]
For example, the complement of $\{2\}$ is indexed by $\{1,3,4\}$; its
$3\times3$ principal submatrix has determinant $1396$ and the sum of its
three diagonal cofactors is $400$, yielding $400/1396=100/349$.

The two marginal error reductions are therefore
\[
 F(\{2\})-F(\{2,3\})=\frac{1761}{19544},\qquad
 F(\{2,4\})-F(\{2,3,4\})=\frac{137}{1518}.
\]
Their difference in the forbidden direction is
\[
 \frac{137}{1518}-\frac{1761}{19544}
 =\frac{2165}{14833896}>0.
\]
Hence selecting index $3$ reduces the error more, not less, after index $4$
has already been selected.  This exactly violates diminishing error
reductions.  Every displayed quantity is an integer or an exact rational.
\end{proof}
\end{cxrefutation}
